\subsection{VAE}
\subsubsection{Introduction to VAE}
The Variational Auto-encoder (VAE) \cite{kingma2014} is an artificial neural network whose encoder and decoder are connected through a probabilistic latent space, corresponding to the parameters of a variational distribution.
\begin{figure}[H]
    \centering
    \includegraphics[width=1\linewidth]{assets/vae_structure.png}
    \caption{Structure of a Variational Auto-encoder, picture from \cite{shende2023}}
    \label{fig:vae structure}
\end{figure}
The encoder receives as input a tensor representing an image extracted from a dataset, and it maps it into a distribution within the latent space, parametrized by a mean vector $\mu$ and a log-variance vector $\log\sigma^2$. The network predicts the logarithm of the variance instead of the standard deviation, so that its output is unconstrained in sign while $\sigma$ remains positive.
\\\\
The decoder, instead, maps from the latent space back to the input space.
\\\\
The VAE's parameters are optimized via gradient descent, however the sampling of the latent variable $z$ introduces stochasticity, making the gradient of the loss with respect to the model parameters non-differentiable. A technique known as Reparameterization Trick is employed to address this problem, which expresses the sampling operation in a way that separates the deterministic and stochastic parts, enabling gradient flow through the deterministic part. So instead of directly sampling $z$ as $z\sim N(\mu,\sigma^2)$, which is non-differentiable with respect to $\mu$ and $\sigma$, it sample an auxiliary variable $\epsilon$ from a standard normal distribution $N(0,I)$, and then it expresses $z$ as $z=\mu+\sigma\cdot\epsilon$, which is now a differentiable function of $\mu$ and $\sigma$.
\\
\subsubsection{Components of VAE}
The two primary components of the VAE are:
\paragraph{Encoder}
\hfill \break
The encoder's architecture consists of four 2D convolution layers, one flatten layer and one dense layer. The input received by this module, which is a 1-channel $28\times28$ image, passes through the convolutional layers -- using 32, 64, 64 and 64 filters respectively -- to extract features, then the feature maps is flattened and its dimensionality is reduced using a dense layer; finally the encoder outputs the mean and the log-variance of the latent space distribution and samples the latent vector $z$.
\begin{table}[H]
    \centering
    \begin{tabular}{l c r}
    \textbf{Layer} & \textbf{Output shape} & \textbf{Parameters} \\
    \hline
    Input & \( 28 \times 28 \times 1 \) & -- \\
    Convolution & \( 28 \times 28 \times 32 \) & 320 \\
    Convolution & \( 14 \times 14 \times 64 \) & 18{,}496 \\
    Convolution & \( 14 \times 14 \times 64 \) & 36{,}928 \\
    Convolution & \( 14 \times 14 \times 64 \) & 36{,}928 \\
    Flatten & \( 12544 \) & -- \\
    Dense & \( 32 \) & 401{,}440 \\
    Dense & \( 20 \) & 660 \\
    Dense & \( 20 \) & 660 \\
    Sampling & \( 20 \) & -- \\
    \hline
    \textbf{Total} & & \textbf{495{,}432} \\
    \end{tabular}
    \caption{Structure of the VAE encoder, shown for a latent dimension of 20. The two dense heads produce the mean and the log-variance of the latent distribution, and the sampling layer applies the reparameterization trick.}
    \label{tab:vae encoder}
\end{table}
\paragraph{Decoder}
\hfill \break
The decoder's architecture consists of a dense layer, a reshape layer and two 2D transposed convolution layers. The input of the module is the latent vector that, through the dense and reshape layers, it is transformed back into a tensor, matching the shape of the last convolution layer of the encoder; finally the convolution layers generate the reconstructed image with the same number of channels as the original input image.
\begin{table}[H]
    \centering
    \begin{tabular}{l c r}
    \textbf{Layer} & \textbf{Output shape} & \textbf{Parameters} \\
    \hline
    Input & \( 20 \) & -- \\
    Dense & \( 12544 \) & 263{,}424 \\
    Reshape & \( 14 \times 14 \times 64 \) & -- \\
    Transposed convolution & \( 28 \times 28 \times 32 \) & 18{,}464 \\
    Transposed convolution & \( 28 \times 28 \times 1 \) & 289 \\
    \hline
    \textbf{Total} & & \textbf{282{,}177} \\
    \end{tabular}
    \caption{Structure of the VAE decoder, shown for a latent dimension of 20. With a latent dimension of 30 the total becomes 407{,}617 parameters.}
    \label{tab:vae decoder}
\end{table}
\subsubsection{Training Process}
The optimization of VAE is achieved through the maximization of the Evidence Lower Bound (ELBO), consisting of two terms: the Reconstruction Loss and the Regularization Term (Kullback--Leibler Divergence).\\Given the input data $x$, the latent variable $z$, the encoder network $q_\phi(z\vert x)$, the decoder network $p_\theta(x\vert z)$, the prior distribution over the latent variables $p(z)$, the reconstruction loss $\mathbb{E}_{q_{\phi}(z\vert x)}[\log p_\theta(x\vert z)]$, and the KL divergence $\mathrm{KL}(q_\phi(z\vert x)\ \Vert\ p(z))$, the loss is computed as:
\\
\begin{equation}
L(\theta,\phi;x)=\mathbb{E}_{q_{\phi}(z\vert x)}[\log p_\theta(x\vert z)]-\mathrm{KL}(q_\phi(z\vert x)\ \Vert\ p(z))
\end{equation}
The VAE is trained with a custom loop, consisting of three main steps: forward pass, loss calculation and backward pass. During the forward pass, the input data is passed through the encoder to get the mean and log variance of the latent distribution, then the latent vector $z$ is sampled using the reparameterization trick, and finally z is passed through the decoder to get the reconstructed data. The loss calculation step computes the reconstruction loss between the input data and the reconstructed data, the KL divergence, and then sums them up to compute the total loss. Finally, in the backward pass the total loss is backpropagated to update the parameters of the encoder and decoder.
\\\\
The model is trained for up to 100 epochs using a batch size of 100 and an Adam optimizer \cite{kingma2015adam} with learning rate of 0.001, with early stopping monitoring the validation loss.
\\\\
\textbf{KL warm-up.} Training the model as described above turned out to be unreliable: repeating it with different random seeds produced generators of very different quality, and the weaker runs showed partial posterior collapse, a known failure mode in which the encoder stops using part of the latent space and the decoder learns to ignore it. To counter it we scale the KL term by a coefficient that grows linearly from 0 to 1 over the first ten epochs. Early in training the latent code can become informative without being pulled back towards the prior, and the regularization reaches full strength only afterwards. The evaluation always uses a coefficient of 1, so validation losses remain comparable across epochs and across runs.
\\
Without warm-up the average KL divergence at convergence was between 4 and 7 nats and the quality of the resulting prior varied by a factor of 3.5 across seeds; with warm-up it is between 16 and 23 nats and every model improves, the best representation error at \( k = 30 \) falling from 0.0103 to 0.0062. Across the two seeds of each configuration the remaining spread is a factor of 2.5, 1.3 and 1.2; that figure and the 3.5 above are measured on different images and over a different number of seeds, so they are given side by side rather than as one ratio that shrank. The rule used to choose between runs was fixed before they were carried out: two seeds per configuration, and the one with the lower validation loss is kept.
\\\\
Two latent dimensions are used throughout. The first, \( k = 20 \), is the value \cite{bora2017} use on MNIST, so that the comparison with their result is direct. The second, \( k = 30 \), tests a suggestion they leave open, that a representation dimension above twenty would make better use of additional measurements. Section 6.5 reports what the experiment says about it.
\subsubsection{The architecture of the reference paper}

The convolutional networks described above are our own design. The experiments of Bora et al. \cite{bora2017} on MNIST instead use a fully connected auto-encoder, with a recognition network of shape \( 784-500-500-k \) and a generator of shape \( k-500-500-784 \), a latent dimension of 20, and the same optimizer, batch size and learning rate we use. Since our generator differs from theirs, a difference between our results and the numbers they publish could be attributed either to the method or to the model.
\\
To separate the two, we also trained their architecture, and we report it as a separate prior in the Results section. The paper specifies the layer widths but not the activation function; we use softplus, following the original variational auto-encoder \cite{kingma2014} the architecture is taken from.

\subsection{Model results}
The following graph shows the output mean vector produced by the encoder with the test set provided by the MNIST dataset. For the sake of simplicity, this result is obtained with a latent dimension equal to 2.
\begin{figure}[H]
    \centering
    \includegraphics[width=0.89\linewidth]{assets/output_encoder.png}
    \caption{Plotted output of the encoder, each color represents a specific digit}
    \label{fig:vae latent space}
\end{figure}
\hfill \break
By observing the data in the graph, we can predict the behaviour of the decoder. Infact, if we assume it receives as input the vector $[-2,-1]$, we expect the output to depict with good precision the number 0, as represented by the blue color in the lateral colorbar. The image reconstruction is shown in the image below.
\begin{figure}[H]
    \centering
    \includegraphics[width=0.5\linewidth]{assets/output_decoder.png}
    \caption{Output of the decoder for the input vector $[-2,-1]$}
    \label{fig:vae decoder output}
\end{figure}